\section{Related Work}
\label{sec:related}

\subsection{Panoramic layout estimation and model-assisted annotation}

Panoramic room-layout estimation converts an equirectangular image into a compact description of enclosing surfaces and their boundaries. HorizonNet predicts ceiling--wall, floor--wall, and wall--wall signals as one-dimensional sequences and reconstructs a layout through geometry-aware post-processing \cite{sun2019horizonnet}. HoHoNet retains the horizontal inductive bias while broadening the representation to layout, depth, and semantic outputs \cite{sun2021hohonet}; LGT-Net more recently uses geometry-aware transformer structure for the single-panorama layout problem \cite{jiang2022lgtnet}. These models make an initial proposal inexpensive, but their benchmark accuracy does not remove the need to detect topological errors, repair uncertain corners, or judge cases outside their training distribution. ZInD further demonstrates that real homes contain a broad layout distribution and that building a large geometric resource requires substantial annotation effort \cite{cruz2021zind}.

Human--machine annotation research provides several mechanisms for reducing this effort. Lean Crowdsourcing combines a model of worker skill, item difficulty, and computer-vision estimates to collect a variable number of annotations \cite{branson2017lean}. Its multiclass extension models worker behaviour and online risk to reduce redundancy and expert review \cite{vanhorn2018leanmulticlass}. More recent work examines practical choices for model-assisted image labelling, including learned representations, simulated human queries, and effort--accuracy trade-offs \cite{liao2021goodpractices}. Together, these studies establish the value of adaptive evidence collection and machine assistance.

Our problem differs in two respects. The annotation is a structured geometric object rather than a class label, so a delivered output can be numerically similar yet structurally unusable. In addition, the system chooses among eligible workers for a task after an independent calibration stage; it is not primarily choosing how many labels to collect or when to stop. We therefore treat model assistance as one experimental factor and focus the principal contribution on worker profiling and routing under explicit evidence support.

\subsection{Worker characterization under sparse and heterogeneous evidence}

The literature on repeated annotation begins with models that infer latent labels and observer-specific error rates. The Dawid--Skene model is foundational because it separates the object being estimated from systematic observer error \cite{dawid1979observer}. Later work relaxes homogeneous-worker assumptions. Yan et al. model expertise as dependent on the instance, showing that a worker may know different parts of the input space \cite{yan2010expertise}. Jin et al. separate difficulty from subjectivity and represent worker preferences when multiple responses can be defensible \cite{jin2018subjectivity}. Burrell and Schoenebeck test common assumptions directly on crowd datasets and document forms of worker and task heterogeneity that a constant accuracy probability cannot express \cite{burrell2023conventional}.

These contributions caution against both over-aggregation and over-interpretation. Agreement can indicate a stable solution in a geometrically ambiguous task, but it is not automatically evidence of correctness. Conversely, disagreement can arise from a valid alternative layout rather than low worker quality. Reference-based evidence also has limits: an operational reference is a frozen scoring target, not a declaration that all room geometry is uniquely observable from one panorama.

The present characterization is intentionally small. It does not estimate an arbitrary high-dimensional ``scene--worker reliability'' surface. Instead, it separates three outcomes with distinct audit meanings. Reference-based accuracy measures proximity to the frozen operational reference. Stable peer compatibility measures whether a worker's output lies near supported peer structure after multimodal cases are excluded. Worker-attributable structural invalidity measures failures that make the output unusable and for which the attribution record supports a worker cause. Conditional effects are limited to predeclared risk or failure families and are retained only when an independent calibration stage provides enough support. This design adapts established ideas---agreement measures, mixed models, and shrinkage---to a setting where sparse evidence can otherwise make detailed profiles look more certain than they are.

\subsection{Adaptive assignment, support, and prospective policy evaluation}

Adaptive task assignment is an established research direction. Ho et al. jointly study assignment and label inference for heterogeneous classification tasks and show why worker diversity can make adaptive allocation useful \cite{ho2013adaptive}. Kang and Tay formulate task recommendation using category-specific worker preferences and reliabilities learned over time with gold tasks \cite{kang2018recommendation}. Online allocation with reusable agents further shows that rejection and resource return are part of the assignment problem, rather than incidental implementation details \cite{sumita2022reusable}. These studies directly preclude a broad novelty claim for ``adaptive worker assignment'' or ``task-related expertise.'' Our contribution instead concerns the evidence boundary that permits a conditional adjustment in a structured annotation deployment.

Support is central because a routing score is a policy, not merely a descriptive model. If calibration contains no comparable task for a worker--risk combination, the direction of a local adjustment is not empirically anchored. Similar problems arise in offline decision-making. SPIBB constrains improvement by reverting toward a baseline in uncertain regions \cite{laroche2019spibb}. Work on deficient-support contextual bandits shows how off-policy learning can fail when actions have zero probability under the logging policy and evaluates restrictions on action or policy spaces \cite{sachdeva2020support}. These adjacent methods do not solve panoramic annotation routing, but they support the general deployment principle that unsupported flexibility should be constrained.

Prospective evaluation is also necessary. Historical replay can test implementation, estimate feasibility, or expose a policy that almost never differs from its baseline. It cannot establish the effect of a policy whose decisions alter which worker sees which task, particularly if historical coverage is uneven. We therefore separate calibration from outcomes, freeze both routing policies before the main experiment, randomize tasks or blocks 50:50, and keep operational constraints symmetric. The comparison asks whether the conditional policy first clears safety gates and only then improves delivery-adjusted quality or resource use. This ordered test distinguishes policy value from the descriptive fit of the underlying worker profile.

Across the three literatures, the unresolved link is therefore specific: structured panoramic outcomes must be characterized in interpretable components; conditional components must earn support in an independent stage; and the resulting routing decision must be tested prospectively against a global policy that remains available as fallback. Sections~\ref{sec:setup} and \ref{sec:profiling} formalize this link.
